\chapter*{Roadmap of This Book}
\addcontentsline{toc}{chapter}{Roadmap of This Book}
\markboth{Roadmap of This Book}{}

This book, \textit{Bitcoin: The Grid Resource}, is written for two audiences
at once: energy professionals who want to understand what Bitcoin mining
actually is as a load, and Bitcoin miners who want to understand the power
market they operate in. Its territory is ERCOT, the Texas grid. If you run a mining site, the book
answers when to slow your machines and when to pause them, which contracts
keep that choice worth having, whether reserves and avoiding the summer transmission
peaks still pay, and what all of it would have earned one full-scale Texas
site, year by year. If you work in power, it answers the same questions from
the grid's side of the meter.

One convention of spelling runs through every chapter: Bitcoin with a capital B is the network, its protocol and the industry built on it; bitcoin in lowercase is the currency and its unit.

The book moves in four parts.

\section*{Part I: Two Systems}

\begin{itemize}
\item \textbf{Chapter~1. Introduction: Bitcoin Mining and Energy.} What mining is, why
it consumes energy at industrial scale, and why electricity is the dominant
operating cost. Includes a look inside a modern ASIC miner and a
range for the network's global consumption, with the reason it cannot be
measured.

\item \textbf{Chapter~2. How Power Systems Work and Why Balance Matters.}
How the grid stays balanced in real time: generation, transmission,
frequency, and what happens when balance is lost, told through the 2021
Winter Storm Uri emergency in ERCOT and the 2003 Northeast Blackout.

\item \textbf{Chapter~3. ERCOT: The Energy-Only Grid.} Why Texas runs
an island grid with no capacity market, how day-ahead and real-time markets
set prices, what scarcity pricing means, and where flexible loads fit in
the market's design.
\end{itemize}

\section*{Part II: The Economics of Mining as a Load}

\begin{itemize}
\item \textbf{Chapter~4. Mining Economics: Hashprice, Efficiency,
Breakeven.} The unit economics of mining, and the single formula
that converts a miner's revenue into the language of power markets:
the breakeven price in dollars per megawatt-hour.

\item \textbf{Chapter~5. The Flexible Load: Curtailment and Power
Optimization.} What price-adaptive operation earned in the chapter's backtests across all eight ERCOT load zones. When to run, when to overclock,
when to pause, and why optimization matters most when margins are thin.
Its answer is a price-response loop, three rules and no forecast, which in ERCOT's West load zone captures 93 to 98
percent of what perfect foresight would gain; Chapter~9 puts the dollars on it.
\end{itemize}

\section*{Part III: Monetizing Flexibility in ERCOT}

\begin{itemize}
\item \textbf{Chapter~6. Buying Power: Contracts, Price Arbitrage, and
the Strike Price.} How miners actually buy electricity, from spot
indexes to fixed-price contracts and standardized blocks; and the
real-time strike price, the backstop that steps a fleet out of price
spikes, valued against real ERCOT price distributions.

\item \textbf{Chapter~7. Ancillary Services: Getting Paid to Stand By.}
The products ERCOT buys to keep the grid stable, what they paid from
the 2021 crisis through today, and why those prices collapsed.

\item \textbf{Chapter~8. 4CP, Demand Response, and the Programs Beyond
Price.} Transmission-charge avoidance and the 12CP rule proposed to
replace 4CP, demand-response programs, and what curtailment is paid for
outside the energy market.

\item \textbf{Chapter~9. The Flexibility Stack.}
The part's capstone: one full-scale mining site runs everything this
book has priced (the strike, the power target, reserves, and 4CP
avoidance) under the contracts of Chapter~6, interactions included. The
chapter sets the site's annual switches as an operator could have set
them in advance, and prints the best annual choice in hindsight beside
it as a ceiling. It then compares that site with two alternatives: the same
site never using its flexibility, and a mine behind the meter that buys
cheap power and gives most of its flexibility away. In the chapter's backtest, against the modeled site
that never used its flexibility, the stack an operator could have
planned was worth 8.0 to 29.5 percent of that site's margin after
power and grid charges, year by
year, and \$123~million over 2021 to 2025. Nearly half of it came in
one week of Winter Storm Uri.
\end{itemize}

\section*{Part IV: Outlook}

\begin{itemize}
\item \textbf{Chapter~10. What Comes Next.} Three sums you can
check: the flexibility stack as hashprice keeps falling, the halving
schedule against any price the world economy can bear, and the security
budget and who pays it when the subsidy is gone, with the grid, the fee
market and the large holder as the candidates. It closes on the case that mining
becomes a standard grid resource because Bitcoin's halving schedule makes it one.
\end{itemize}

The appendix documents the methodology and the companion repository: every
number this book computed can be regenerated from the published analysis
code, and every number it did not compute carries a citation.
